%% ════════════════════════════════════════════════════════════════════════════
\section{Introduction}
\label{sec:intro}

Ambient air pollution is among the most pressing environmental health challenges in modern India, contributing to an estimated 1.67 million premature deaths in 2019 \cite{pandey2021} and imposing substantial socio-economic burdens. While national monitoring efforts have expanded dramatically under the National Clean Air Programme (NCAP), reliable predictive capabilities remain vital. Robust short-range forecasts of the daily Air Quality Index (AQI) afford municipal authorities, environmental regulators, and vulnerable populations actionable lead time to implement emergency interventions, regulate industrial emissions, and restrict vehicular circulation.

Gujarat, located along India's western industrial corridor, represents a critical yet under-studied theatre for atmospheric modeling. Hosting vast petrochemical complexes, chemical manufacturing zones, and rapidly expanding urban centres (notably Ahmedabad, Surat, Ankleshwar, and Vapi), Gujarat regularly records extreme concentrations of respirable particulate matter ($PM_{2.5}, PM_{10}$) and toxic gaseous pollutants ($SO_2, NO_2, CO$). In contrast to the National Capital Region (NCR) around Delhi, which has received extensive analytical attention, empirical machine learning research focused specifically on Gujarat's distinct geographical and industrial micro-climates remains sparse.

Over the past decade, a surge of applied machine learning literature has investigated urban AQI prediction across Indian municipalities. Authors have deployed linear estimators, support vector machines, Random Forests, Gradient Boosted Decision Trees (XGBoost, LightGBM, CatBoost), and diverse deep learning architectures (e.g., LSTM, GRU). These publications routinely report coefficients of determination ($R^2$) exceeding $0.80$ to $0.99$, concluding that nonlinear machine learning has effectively solved urban air quality forecasting \cite{maltare2023}\cite{sidhu2024}\cite{mishra2025}\cite{jayapradha2025}\cite{chauhan2026}.

In this paper, we argue that these reported results are heavily inflated by methodological flaws in experimental design. Specifically, two systematic modes of data leakage pervade published studies:

\paragraph{Concurrent Input Leakage (Nowcasting vs. Forecasting Conflation)}
The Central Pollution Control Board (CPCB) prescribes the official daily AQI as the non-linear algebraic supremum across criteria pollutant sub-indices:
\begin{equation}
  \text{AQI}_t = \max_{p \in \mathcal{P}} I_p(C_{p,t})
  \label{eq:cpcb_aqi}
\end{equation}
where $\mathcal{P} = \{PM_{2.5}, PM_{10}, NO_2, SO_2, CO, O_3, NH_3\}$ is the set of monitored criteria pollutants, $C_{p,t}$ is the daily mean concentration of pollutant $p$ on calendar day $t$ (a 24-hour mean for most pollutants; CPCB uses an 8-hour rolling average for CO and O$_3$), and $I_p(C_{p,t})$ is the piecewise-linear sub-index conversion function specified by Indian regulatory guidelines. Numerous applied studies train models where concurrent pollutant concentrations $\{C_{p,t}\}_{p \in \mathcal{P}}$ serve as independent input features to predict that same day's aggregate index $\text{AQI}_t$ \cite{sidhu2024}\cite{mishra2025}\cite{jayapradha2025}\cite{chauhan2026}. Because same-day pollutant concentrations cannot be observed hours or days in advance, this configuration does not represent operational forecasting; it performs static algebraic reconstruction or nowcasting. (In our own Setup~A, concurrent inputs reconstruct the target with $R^2 = 0.815$ rather than $1.0$, likely because the target is the official CPCB daily bulletin index whereas the pollutant inputs come from a separate station-level Kaggle archive.)

\paragraph{Random-Split Autocorrelation Leakage}
Urban daily AQI series exhibit intense temporal persistence and first-order autocorrelation ($\rho_1 \approx 0.75$ in Ahmedabad, above $0.83$ in Ankleshwar and Gandhinagar, and $0.90$ in Vapi).\footnote{Computed directly from the processed Parquet files via \texttt{pandas.Series.autocorr(lag=1)} on valid (non-NaN) AQI observations: Ahmedabad 0.753, Ankleshwar 0.836, Chennai 0.704, Gandhinagar 0.834, Surat 0.710, Vapi 0.901.} When datasets are partitioned using random train-test splitting (e.g., standard random 80/20 cross-validation), adjacent temporal observations ($t-1$ and $t+1$) invariably enter the training partition while day $t$ is placed in the test partition. Because adjacent days share nearly identical meteorological and atmospheric states, the model interpolates between proximal training points rather than projecting forward in time. This provides artificial predictive cues that collapse completely in deployment.

Although the pitfalls of evaluating time-series predictors with shuffled cross-validation have been examined in the forecasting literature \cite{bergmeir2012}\cite{cerqueira2020}, and leakage in predictive modelling has been formalised more broadly \cite{kaufman2012}, their practical effect on Indian AQI studies has not been systematically audited. Furthermore, published studies rarely evaluate simple heuristic benchmarks -- such as the naive Persistence forecaster ($\hat{y}_{t+1} = y_t$) -- or conduct formal statistical hypothesis testing (e.g., the Diebold-Mariano test \cite{diebold1995}) to establish whether machine learning models provide statistically verifiable skill over naive inertia.

\paragraph{Contributions.}
To address these deficiencies and establish rigorous benchmarks, this study delivers three primary contributions:
\begin{enumerate}
  \item \textbf{Empirical Quantification of Leakage Inflation:} We design three nested, controlled evaluation setups (Setups~A, B, and C) applied to municipal series, measuring the inflation associated with concurrent inputs, random train-test partitioning, and temporal distribution shifts.
  \item \textbf{Leak-Free Multi-City Benchmark for Gujarat:} To our knowledge, this is the first leak-free, persistence-benchmarked multi-city evaluation for Gujarat. We implement a strict rolling-origin expanding-window evaluation across five major Gujarat industrial and urban centres (Ahmedabad, Gandhinagar, Ankleshwar, Vapi, Surat) alongside Chennai as a maritime comparative anchor. We evaluate next-day ($h=1$) and three-day ($h=3$) lead times, pairing error distributions with 95\% moving block-bootstrap confidence intervals and Diebold-Mariano significance testing against Persistence.
  \item \textbf{Complete Open-Source Reproducibility Framework:} We release a fully modularised, tested Python package, verified dataset manifests with SHA-256 integrity checksums, and the 138 primary out-of-sample prediction trajectories (supplemented by classical baseline and sensitivity scripts), allowing the entire empirical pipeline to be executed via a single command.
\end{enumerate}

The remainder of this article is organized as follows: Section~\ref{sec:related} reviews related work and methodological pitfalls. Section~\ref{sec:data} details data curation, sensor range validation, and city qualification. Section~\ref{sec:method} formalises the evaluation protocols, leakage guards, candidate forecasters, and hypothesis testing framework. Section~\ref{sec:results} presents the empirical results. Section~\ref{sec:discussion} discusses meteorological mechanisms, practical implications, and limitations. Section~\ref{sec:conclusion} concludes.

